\scanpage{101}
\mspage{3.6}
Corollary 3.9: Let $R\subset S$ be a local algebra of the integrally closed local rings such that the quotients field extension $k\subset K$ is galois. Then $r(S/R)=$ the number of minimal primes $\underline p$ in $S^f$ such that $\underline p\cap R=0=$ the order of the inertia group of $S$ over $R$.

\uline{Proof:} Let $U$ be the inertial ring of $S$ over $R$. Since $U$ is an unramified local $R$-algebra, we have by lemma 2.1 that $S_R^f\approx S_U^f$. However $U$ is a transporter for $S$ over $R$ and hence $S_R^f\approx S_U^f=S\otimes_U S$. For a prime ideal $\underline p$ in $S_R^f$ it is clear that $\underline p\cap R=0$ if and only if $\underline p\cap U=0$, since $U$ is a localization of an integral extension of $R$. Therefore the minimal prime ideals $\underline p$ in $S_R^f$ with $\underline p\cap R=0$ are in one-one correspondence with the prime ideals in $(S\otimes_U S)\otimes_U K_I=K\otimes_{K_I}K$ where $K_I$ is the field of quotients of $U$, i.e. the inertial field of $S$ over $R$. Now $K$ is a galois extension of $K_I$ with $G_I$ the inertia group of $S$ over $R$ as its galois group. We recall that there is one-one correspondence between the prime ideals in $K\otimes_{K_I}K$ and the elements in $G_I$ under $\underline P(\sigma)=\operatorname{Ker}(\widetilde\sigma=K\otimes_{K_I}K\longrightarrow K)$ where $\widetilde\sigma(x\otimes y)=x\sigma(y)$. Therefore the prime ideals $\underline p$ in $S_R^f\approx S\otimes_U S$ with $\underline p\cap R=0$ are in one-one correspondence with the elements in $G_I$ under $\underline p(\sigma)=\underline P(\sigma)\cap S_R^f=\underline P(\sigma)\cap(S\otimes_U S)$. Since $G_I$ is the inertia group of $S$ over $R$, we have in particular that $\sigma(S)=S$ for all $\sigma\in G_I$, and hence $\underline p(\sigma)=\underline P(\sigma)\cap(S\otimes_U S)=\operatorname{Ker}(\widetilde\sigma:S\otimes_U S\longrightarrow S)$ where $\widetilde\sigma(x\otimes y)=x\sigma(y)$. Therefore $S_R^f/\underline p(\sigma)=S\otimes_U S/\underline p(\sigma)=S$ and hence $e(S_R^f/\underline p(\sigma)\mid S)=1$ for all $\sigma\in G_I$. Therefore $r(S/R)=\sum_{\in G_I}e(S^f/\underline p(\sigma)\mid S)=$ the number of prime ideals $\underline p$ in $S_R^f$ with $\underline p\cap R=0\unsure{=|\equiv}$ the order of $G_I$.

Proposition 3.10: Let $R$ be an integrally closed local domain with quotient field $k$. Let $K$ be a finite algebraic extension of $k$ and $S$ the
\mspage{100}
